\subsection*{Appendix E: The Capacity Step-Up Problem}

A technical cohort offers 10 expert review minutes per learner each week, plus a fixed 5-hour live workshop. At 15 learners, review requires 150 minutes, or 2.5 hours, and live teaching raises the direct weekly activity to 7.5 hours, excluding preparation and support. At 30 learners, review requires 5 hours and direct activity reaches 10 hours. If the provider has only 8 sustainable hours per week for this specific program, the larger cohort violates capacity even though the live workshop itself can serve both groups.

The founder could shorten reviews, automate mechanical checks, hire another qualified reviewer, add a 2nd cohort, or change scope and price. These are different business models with different quality and cost implications. Treating all 30 buyers as having the same marginal cost as the 15th ignores the actual scarce resource.

\subsection*{Appendix F: How to Conduct a Unit Economics Review Meeting}

Bring finance, sales, delivery, and customer success perspectives together: even if the founder currently performs all 4 roles. First inspect the customer job and whether the promise is being delivered. Then reconcile paid sales, realized revenue, refunds, and cash receipts. Compare variable delivery with forecast, including support and exceptions. Review acquisition by source, not as 1 averaged CAC. Identify current capacity and the next staffing step. Inspect accounts receivable, payables, and obligations to already-paid customers. Finally choose 1 experiment to improve the binding constraint.

A good review ends with a \textbf{decision}, not merely dashboards: stop an expensive channel, narrow the package, raise the price for a clearly differentiated tier, improve assessment automation, change payment milestones, or decline a project that would monopolize critical time. Record what evidence would reverse the decision next month.

\subsection*{Appendix G: Business-Model Design Under a Full-Time Founder Constraint}

Imagine an early founder with 20 hours per week available after their primary employment obligations. Allocate an illustrative:
\begin{itemize}
  \item 8 hours to customer discovery, marketing, and sales,
  \item 6 hours to direct teaching and feedback delivery,
  \item 4 hours to curriculum and platform asset development,
  \item 2 hours to administration, invoicing, and quality review.
\end{itemize}
That is already 20 hours, without contingency. A paid contract requiring 20 additional delivery hours weekly cannot be treated as free revenue; it requires removing another activity, hiring, or changing scope.

This is 1 reason comparing \textbf{contribution per founder-hour} and \textbf{capability accumulated per hour} is more useful than revenue alone at the beginning. A project generating ₹2 lakh of positive contribution but exhausting all available time for 3 months may block a smaller program that establishes repeatable paid demand. The preferred choice is a strategic decision, but the schedule should expose the trade-off honestly.

\subsection*{Appendix H: A Final Accuracy Checklist for Financial Claims}

Before publishing a revenue, profit, margin, LTV, or customer-impact claim, ask:
\begin{enumerate}
  \item What is the unit of analysis?
  \item Which exact calendar period is covered?
  \item Are amounts reported gross or net of trade discounts and taxes?
  \item Are amounts booked, invoiced, collected, or recognized?
  \item Which delivery, support, and acquisition costs are included?
  \item What are the remaining fixed commitments?
  \item Have refunds, concessions, and failed collections been accounted for?
  \item Does the figure assume uncompensated founder labor?
  \item Which future retention observations actually exist in historical cohorts?
  \item Are prepaid service obligations still outstanding?
  \item Is the measure consistent with the underlying books and source contracts?
  \item What does the figure \emph{not} establish?
\end{enumerate}

A serious founder is not required to produce audited accounts for every pilot, but should be able to explain the construction and limits of every number they use to make a consequential decision. Investors, students, partners, employees, and the founder's own household ultimately bear the costs of numbers that are impressive but wrong.
